% ---------------------------------------------------------------------------
% Related Work draft for the DIF-Gaussian extension paper.
% Assumes \method{} is defined in introduction.tex (include both files in the
% main document). Citation keys are placeholders; suggested references are
% listed in the comment block at the end. VERIFY all bibliographic details.
% ---------------------------------------------------------------------------

\section{Related Work}
\label{sec:related}

\subsection{Analytic and Model-Based Iterative Reconstruction}
\label{sec:related:analytic}

Analytic reconstruction inverts the imaging model in closed form. FDK
\cite{feldkamp1984practical} remains the clinical workhorse because it is fast
and deterministic, but its accuracy rests on dense angular sampling: under
extremely sparse-view acquisition it suffers from view-starvation artifacts,
including streaks, angular aliasing, and eroded boundaries. Model-based
iterative reconstruction casts the problem as regularized optimization of a
data-fidelity term, and compressed-sensing formulations with transform-sparsity
priors such as total variation substantially suppress these artifacts
\cite{sidky2008image}, at the cost of minutes-to-hours of per-patient
optimization, sensitivity to regularization hyperparameters, and dependence on
an accurate system model. In contrast, our method reconstructs in a single
feed-forward pass and involves no per-scan optimization.

\subsection{Deep Learning--Based Sparse-View Reconstruction}
\label{sec:related:dl}

Supervised deep learning replaces the inverse operator with a learned network.
Most approaches denoise FDK volumes with 2D convolutional networks applied
slice-by-slice \cite{chen2017redcnn,kang2017deep}, which leaves inter-slice
inconsistency and residual 3D streaks; 3D variants improve coherence
\cite{shan2019competitive} but require dense-view training data and volumetric
memory. Learned iterative and dual-domain methods
\cite{adler2018learned,peng2022dudotrans} incorporate more imaging structure
but remain tied to fixed 2D processing pipelines with considerable inference
cost. In contrast, our method reasons directly in 3D---features are lifted
onto 3D anchor tokens and fused across views under the actual projection
geometry---and generalizes to unseen patients and view counts.

\subsection{Implicit Neural Representations for CT}
\label{sec:related:inr}

Implicit neural representations model the volume as a continuous function
optimized per case. INR-based attenuation fields such as NAF, NeAT, and
Semi-INR \cite{zang2021naf,cai2023neat,cai2024semiinr} reconstruct sparse-view
CBCT with strong view consistency, and radiative Gaussian splatting methods
such as R2-Gaussian and X-Gaussian \cite{zhang2024r2gaussian,cai2024xgaussian}
accelerate this with explicit 3D Gaussian primitives. All of them, however,
optimize every new patient from scratch, take tens of minutes per case, and
learn no cross-patient prior. In contrast, our method distills a similar
continuous Gaussian representation into a feed-forward encoder trained across
patients and reconstructs a new case in a single forward pass.

\subsection{Generalizable Implicit-Field Reconstruction}
\label{sec:related:generalizable}

Closest to our work are feed-forward implicit-field methods that learn a
cross-patient prior for sparse-view reconstruction. DIF-Gaussian
\cite{lin2024dif} maps projections to anisotropic 3D Gaussians carrying deep
implicit features, decoded continuously by Gaussian-weighted interpolation,
and generalizes across patients and view counts with bounded memory;
feed-forward image-to-volume approaches such as X2CT-GAN and FBI
\cite{ying2019x2ct,qiu2023fbi} share this single-pass spirit with weaker
geometric grounding. However, the way these methods aggregate cross-view
evidence remains view-agnostic---typically channel-wise maximum pooling on the
coarsest bottleneck scale---which is geometry-blind, hard-selects rather than
reconciles conflicting observations, and leaves anchors spatially isolated. In
contrast, we retain the Gaussian implicit field but replace view-agnostic
fusion with a geometry-aware multiview transformer that conditions cross-view
attention on explicit ray geometry, multi-scale sampling, and FOV validity,
followed by local 3D self-attention across anchors.

\subsection{Geometry-Aware Attention and Multi-View Feature Fusion}
\label{sec:related:attention}

Geometry conditioning has proven effective when fusing information across
views. Multi-view stereo networks couple feature matching with camera
geometry through differentiable homography warping \cite{yao2018mvsnet}, and
epipolar or cross-attention transformers propagate features along
geometrically corresponding image locations for novel-view synthesis
\cite{ravi2020epipolar,yu2021pixelnerf,wang2021ibrnet}. These mechanisms,
however, were designed for per-scene neural rendering or 2D correspondence
search: attention lives in image space, and geometry enters only as
image-to-image correspondences. In contrast, our attention operates on 3D
anchor tokens inside a generalizable reconstruction network: each anchor
attends to ray-encoded, FOV-masked observations from all views and scales,
and neighboring anchors exchange information through windowed 3D
self-attention with physical relative-position biases.

% ---------------------------------------------------------------------------
% Additional suggested references (VERIFY details):
%
%  chen2017redcnn     - H. Chen et al., "Low-dose CT with a residual
%      encoder-decoder convolutional neural network", IEEE TMI, 2017.
%  shan2019competitive - H. Shan et al., "Competitive performance of a
%      modularized deep neural network compared to commercial algorithms for
%      low-dose CT image reconstruction", Nature Machine Intelligence, 2019.
%  peng2022dudotrans  - C. Peng et al., "DuDoTrans: Dual-domain transformer
%      provides high quality and robust learning for sparse-view CT image
%      reconstruction", ECCV 2022.
%  ying2019x2ct       - X. Ying et al., "X2CT-GAN: Reconstructing CT from
%      biplanar X-rays with generative adversarial networks", CVPR 2019.
%  qiu2023fbi         - "FBI: Fourier basis inpainting for sparse-view CT"
%      (verify authors/venue/year).
%  yao2018mvsnet      - Y. Yao et al., "MVSNet: Depth inference for
%      unstructured multi-view stereo", ECCV 2018.
%  ravi2020epipolar   - "Epipolar transformers" (CVPR 2020; verify authors).
%  yu2021pixelnerf    - A. Yu et al., "pixelNeRF: Neural radiance fields from
%      one or few images", CVPR 2021.
%  wang2021ibrnet     - Q. Wang et al., "IBRNet: Learning multi-view image
%     -based rendering", CVPR 2021.
% ---------------------------------------------------------------------------
